% Chapter 11 -- Claim: Value-at-Risk Models Measure Risk Reliably
% Requires: amsmath, amssymb, booktabs, array, tabularx, graphicx, tcolorbox (breakable, skins), hyperref or url.
% Defines the callout box only if the preamble has not.
\providecolor{fenavy}{HTML}{1B2A6B}
\providecolor{feshade}{HTML}{F3F5FA}
\makeatletter
\@ifundefined{callout}{%
  \newtcolorbox{callout}{breakable, enhanced,
    colback=white, colframe=gray!70, boxrule=0.5pt, arc=1mm,
    borderline west={2.5pt}{0pt}{gray!70},
    left=8pt, right=8pt, top=4pt, bottom=4pt}%
}{}
% Shaded box for optional technical material that general readers may skip.
\@ifundefined{technote}{%
  \newtcolorbox{technote}{breakable, enhanced,
    colback=feshade, colframe=fenavy!60, boxrule=0.5pt, arc=1mm,
    left=8pt, right=8pt, top=4pt, bottom=4pt}%
}{}
\makeatother

\chapter{Claim: Value-at-Risk Models Measure Risk Reliably}
\label{ch:08}
\providecommand{\term}[1]{\emph{#1}}

A weather forecast that gives a 1\% chance of a severe storm tomorrow makes a promise that can be checked: over many days, severe storms should arrive on about one day in a hundred. Banks make the same kind of promise about their losses every day, and regulators check it. This chapter asks whether those promises are kept, especially in the storms that matter most.

In the early 1990s large banks needed a single number that summarized the market risk of a trading book, the stocks, bonds and other securities a bank buys and sells for profit. J.P. Morgan's RiskMetrics system, made public in 1994, popularized one: \textbf{Value at Risk}\index{Value at Risk (VaR)}, the loss that should be exceeded only on a stated fraction of days. A bank that reports a one-day 99\% VaR of \$50 million is claiming that it will lose more than \$50 million on about one trading day in a hundred. In 1996 the \term{Basel Committee} on Banking Supervision, the international group of bank regulators, allowed banks to set their capital for market risk, the money of their own that they must hold to absorb losses, with their own VaR models, provided the models passed a regular backtest.


The question matters beyond banks. Every investor who sets a risk limit, a stop-loss or a leverage ratio from a model is relying on the same kind of forecast. A day on which the loss is larger than the VaR is called a \term{breach}\index{VaR exceedance}. If a bank's 99\% VaR is \$50 million, a breach is a day on which it loses more than \$50 million, and a good model should produce about one such day in every hundred. The timing matters as well as the count. The breaches should be scattered randomly through the years, not bunched together in a few bad weeks, because a model whose breaches all arrive during a crisis fails exactly when it is needed most.
\begin{figure}[htbp]
\centering
\scalebox{0.9}{%
\begin{tikzpicture}
\begin{axis}[
  width=0.75\linewidth, height=0.38\linewidth,
  axis lines=left, axis line style={-stealth},
  xmin=0, xmax=255, ymin=-95, ymax=70,
  xtick={0,50,...,250}, ytick={-50,0,50},
  yticklabels={$-50$,$0$,$50$},
  xlabel={Day},
  ylabel={Daily P\&L (\$ million)},
  clip=false,
]
  % One year of simulated daily P&L, normal with 99% VaR of $50 million
  \addplot[ycomb, fegray!70, no marks] coordinates {%
(1,7.4) (2,17.7) (3,7.1) (4,-28.0) (5,19.5) (6,9.6) (7,-11.5) (8,12.5) (9,7.8) (10,6.3) (11,0.6) (12,11.8) (13,-15.8) (14,-3.5) (15,-10.4) (16,12.9) (17,0.9) (18,-6.3) (19,-16.8) (20,-5.5) (21,0.2) (22,-5.9) (23,27.8) (24,21.6) (26,-40.6) (27,-3.8) (28,-9.1) (29,4.6) (30,4.7) (31,45.5) (32,-23.9) (33,-8.1) (34,43.9) (35,13.9) (36,14.3) (37,-11.0) (38,-35.4) (39,3.6) (40,2.3) (41,-26.4) (42,-14.7) (43,-1.5) (44,-20.3) (45,-2.1) (46,2.1) (47,0.8) (48,-10.9) (49,12.8) (50,19.2) (51,6.9) (52,-17.6) (53,15.7) (54,-10.8) (55,18.9) (56,-23.0) (57,19.7) (58,-0.4) (59,-26.8) (60,-6.7) (61,1.2) (62,5.9) (63,-21.1) (64,-23.8) (65,4.3) (66,-10.0) (67,5.1) (68,16.3) (69,-35.4) (70,5.5) (71,26.3) (72,-6.4) (73,-17.4) (74,16.2) (75,5.4) (76,19.3) (77,-7.4) (78,-31.9) (79,-2.4) (80,-9.6) (81,16.7) (82,4.2) (83,-35.1) (84,-25.7) (85,19.0) (86,14.6) (87,-13.8) (88,-0.0) (89,9.6) (90,10.1) (91,18.8) (92,5.5) (93,-2.0) (94,-5.6) (95,22.7) (96,-48.4) (97,-3.0) (98,0.7) (99,-30.6) (100,7.2) (101,-14.0) (102,18.5) (103,-2.7) (104,14.4) (105,26.2) (106,8.2) (107,-18.8) (108,-32.6) (109,37.7) (110,-2.4) (111,-14.8) (112,3.1) (113,-4.1) (114,18.3) (115,0.7) (116,0.3) (117,-15.4) (118,10.1) (119,-22.2) (120,14.3) (121,32.8) (122,-32.8) (124,13.3) (125,54.8) (126,-21.5) (127,-26.9) (128,12.7) (129,-18.1) (130,-10.9) (131,-7.5) (132,11.4) (133,-8.7) (134,6.0) (135,-3.8) (136,-18.2) (137,-6.9) (138,-20.4) (139,0.1) (140,-24.2) (141,-23.5) (142,31.3) (143,-1.1) (144,-1.2) (145,11.0) (146,-9.0) (147,-4.9) (148,9.1) (149,6.1) (150,-24.9) (151,17.9) (152,-12.7) (153,-22.7) (154,-19.4) (155,-8.4) (156,35.0) (157,-25.3) (158,3.4) (159,-46.0) (160,-0.0) (161,19.3) (162,-5.1) (163,-13.5) (164,5.0) (165,15.1) (166,14.3) (167,42.4) (168,4.5) (169,-12.7) (170,-2.7) (171,-1.6) (172,2.3) (173,-0.6) (174,3.7) (175,-35.9) (176,17.8) (177,-12.4) (178,-25.2) (179,13.7) (180,28.3) (181,10.6) (182,3.5) (183,-20.0) (184,61.7) (185,18.9) (186,-24.5) (187,-16.8) (188,1.9) (189,-33.4) (190,3.6) (191,-9.9) (192,26.4) (193,20.7) (195,0.9) (196,-34.8) (197,23.9) (198,3.6) (199,11.8) (200,-22.9) (201,39.3) (202,43.4) (203,-22.9) (204,8.0) (205,-14.5) (206,-0.5) (207,-27.2) (208,40.1) (209,-20.8) (210,-6.4) (211,10.8) (212,-13.9) (213,-5.1) (214,-12.1) (215,-2.9) (216,-25.2) (217,-9.4) (218,-4.4) (219,-7.2) (220,1.2) (221,-6.3) (222,16.2) (223,-6.9) (224,-2.9) (225,-14.3) (226,-11.3) (227,-27.2) (228,11.2) (229,-24.6) (230,-16.0) (231,7.7) (232,8.7) (233,-8.6) (234,-43.4) (235,9.0) (236,5.6) (237,-30.4) (238,16.6) (239,-15.1) (240,-24.2) (241,2.1) (242,-3.8) (243,4.4) (244,-34.5) (245,39.0) (246,-13.0) (247,-33.1) (248,13.3) (249,-7.6) (250,7.0)};
  % Breaches: losses larger than the VaR
  \addplot[ycomb, feorange, very thick, mark=*, mark size=1.3pt] coordinates {(25,-58.3) (123,-53.0) (194,-58.3)};
  % VaR line
  \draw[fenavy, very thick, dashed] (axis cs:0,-50) -- (axis cs:250,-50)
    node[right, font=\small, align=left] {99\% VaR:\\$-$\$50 million};
  \node[anchor=south west, font=\small\itshape, text=fegray] at (axis description cs:0.03,1.0) {Hypothetical trading book};
  \node[feorange, font=\footnotesize] at (axis cs:125,-84)
    {3 breaches in 250 days, against about 2.5 expected};
\end{axis}
\end{tikzpicture}%
}
\caption{What a 99\% one-day VaR of \$50 million means. Each bar is one day's profit or loss for a hypothetical trading book over one year, simulated from a normal distribution. The VaR line sits where losses are expected to fall below it on about 1 day in 100. Over 250 trading days the loss exceeds the VaR on 3 days, shown in orange, close to the 2.5 a correct model would produce.}
\label{fig:11-var-intro}
\end{figure}


That decision raised a statistical question with real money attached: how can anyone tell whether a VaR model is right? VaR is unusual among the claims in this book because it is a forecast that checks itself. Every day produces a test, because either the loss exceeded the VaR or it did not. To check banks' models from the outside, \term{Kupiec}\index{Kupiec test} (1995) proposed a test of whether breaches occur as often as the model promises, and \term{Christoffersen} (1998) added a test of whether they arrive independently or in clusters, since a model that is right on average but wrong in every crisis is of little use. The crisis of 2007--09, when breaches came in runs, led regulators to replace VaR with \term{expected shortfall (ES)}, the average loss on the worst days, in the Fundamental Review of the Trading Book, Basel's overhaul of the rules for market risk.
Since the Fundamental Review of the Trading Book (2016), banks set market-risk capital with 97.5\% ES, not 99\% VaR. 


The claim on trial\index{claims-on-trial method} is therefore: \textbf{Value-at-Risk models measure risk reliably. Losses exceed a 99\% VaR on about 1\% of trading days, in calm markets and in crises alike.}

The lab in Section~\ref{sec:11-lab} replicates the backtests of Kupiec (1995) and Christoffersen (1998) on six common VaR models, using free daily U.S. market returns since 1930, in order to check whether any of them measures risk reliably and how they fail in crises.


\section{Foundations}
\label{sec:11-foundations}

\subsection{Value at Risk}

In words, VaR answers a simple question: on a bad day of the kind that comes about once in a hundred, how much could we lose? For a confidence level $1 - \alpha$ and a horizon (one day, ten days), VaR is the loss that is exceeded with probability $\alpha$:
\begin{equation}
  P(\text{Loss} > \mathrm{VaR}) = \alpha
  \label{eq:11-var}
\end{equation}
It is a \term{quantile}\index{quantile|textbf}, or percentile, of the loss distribution. A 99\% one-day VaR is the 1st percentile of tomorrow's return, with the sign flipped. If daily returns are normal with mean zero and standard deviation\index{standard deviation} $\sigma$, then $\mathrm{VaR} = z\sigma$, where $z = 1.645$ at 95\%, 2.326 at 99\% and 3.090 at 99.9\%.

\subsection{Expected shortfall (ES)}
If we simply sort the 250 daily profits and losses in Figure~\ref{fig:11-var-intro} by size, most of them cluster near zero and fewer lie far out on either side, so together they form a bell-shaped curve close to the normal distribution (Figure~\ref{fig:11-var-es}). 

\textbf{Expected shortfall (ES)}\index{expected shortfall (ES)|textbf}, also called conditional VaR, is the average loss on the days when the VaR is exceeded:
\begin{equation}
  \mathrm{ES} = E[\text{Loss} \mid \text{Loss} > \mathrm{VaR}]
  \label{eq:11-es}
\end{equation}
For a normal distribution\index{distribution!normal}, $\mathrm{ES} = \sigma\,\phi(z)/\alpha$, where $\phi$ is the normal density. ES always exceeds VaR at the same confidence level. VaR says where the tail starts. ES says how bad the tail is (Figure~\ref{fig:11-var-es}).

\begin{figure}[htbp]
\centering
\scalebox{0.9}{%
\begin{tikzpicture}[declare function={npdf(\x) = exp(-\x*\x/2)/sqrt(2*pi);}]
\begin{axis}[
  width=0.66\linewidth, height=0.36\linewidth,
  axis lines=left, axis line style={-stealth},
  xmin=-6, xmax=4, ymin=0, ymax=0.45,
  xtick={-2.33,0,2}, xticklabels={$-2.33$,$0$,$2$},
  ytick=\empty,
  xlabel={Tomorrow's return, in daily standard deviations},
  ylabel={Likelihood},
  samples=200, clip=false,
]
  \addplot[name path=curve, draw=none, domain=-4.2:4] {npdf(x)};
  \path[name path=base] (axis cs:-4.2,0) -- (axis cs:4,0);
  \addplot[fenavy] fill between[of=curve and base, soft clip={domain=-4.2:-2.33}];
  \addplot[fegray, thick, domain=-4.2:4] {npdf(x)};

  % 99% VaR: label sits on top of its line
  \draw[fenavy, very thick] (axis cs:-2.33,0) -- (axis cs:-2.33,0.30);
  \node[fenavy, anchor=south, align=center, font=\small] at (axis cs:-2.33,0.30)
    {\textbf{99\% VaR = 2.33 s.d.}\\the worst 1\% of days\\start here};

  % 99% ES: label sits to the left of its line
  \draw[thick, dashed] (axis cs:-2.67,0) -- (axis cs:-2.67,0.13);
  \node[anchor=east, align=right, font=\small] at (axis cs:-0.7,0.13)
    {\textbf{99\% ES = 2.67 s.d.}\\average loss\\on those days};
\end{axis}
\end{tikzpicture}%
}
\caption{VaR marks where the worst days begin, and expected shortfall measures how bad they are. VaR and expected shortfall for normal returns. The shaded tail holds the worst 1\% of days. The 99\% VaR, 2.33 standard deviations below zero, marks where that tail begins. The 99\% ES, 2.67 standard deviations, is the average loss on the days in the tail, so it is always larger than the VaR.}
\label{fig:11-var-es}
\end{figure}

\subsection{From one-day to 10-day VaR: the square-root rule and its bias}


A bank cannot always sell a losing position in a day. Large or illiquid positions may take a week or two to unwind, so regulators set capital with a 10-day VaR rather than a one-day VaR. Ten-day losses are hard to estimate directly, because even 20 years of data contain only about 500 non-overlapping 10-day periods. In practice, banks therefore compute a one-day VaR and scale it up. If daily returns are independent with the same variance\index{variance}, the $h$-day VaR is $\sqrt{h}$ times the one-day VaR\index{Value at Risk (VaR)!square-root-of-time rule}, the same square-root rule used to annualize volatility (Section~\ref{sec:10-foundations}).

The rule assumes that today's volatility will last for the whole horizon, but volatility mean-reverts (Chapter~\ref{ch:07}). After a shock, volatility is likely to fall, so the rule overstates the 10-day risk. After a calm period, volatility is likely to rise, so the rule understates it. The second error is the dangerous one. A 10-day VaR scaled from a quiet year looks lowest just before risk returns, as it did for many banks in early 2007.

\textbf{Worked example.} A \$10 million position in a stock index has a daily volatility of 1.2\%. Its normal 99\% one-day VaR is $2.326 \times 1.2\% \times \$10\text{m} \approx \$279{,}000$, and its 99\% ES is $2.665 \times 1.2\% \times \$10\text{m} \approx \$320{,}000$. Its 97.5\% ES is about \$281{,}000, almost the same as the 99\% VaR, which is why regulators chose 97.5\% when they moved from VaR to ES: the new measure kept capital at about the same level while capturing the size of tail losses. By the square-root rule, its 10-day 99\% VaR is $\sqrt{10} \times \$279{,}000 \approx \$883{,}000$.

\subsection{When diversification raises VaR}

Diversification\index{diversification} should never make a portfolio riskier, so the risk of two positions held together should be no more than the sum of their separate risks. A risk measure with this property is called \term{subadditive}\index{subadditivity|textbf}\index{coherent risk measure} (Artzner, Delbaen, Eber and Heath 1999). ES always has it, but VaR does not.

Here is an example. Two bonds each have a 4\% chance of default, and each loses 100 if it defaults. If both are held, the chance that at least one defaults is $1 - 0.96^2 = 7.84\%$.

\begingroup
\small
\setlength{\LTpre}{6pt}\setlength{\LTpost}{6pt}
\begin{longtable}{@{}>{\raggedright\arraybackslash}p{0.28\linewidth}>{\raggedright\arraybackslash}p{0.18\linewidth}>{\raggedright\arraybackslash}p{\dimexpr 0.9\linewidth-0.28\linewidth-0.18\linewidth-4\tabcolsep\relax}@{}}
\toprule
Position & Chance of a loss & 95\% VaR \\
\midrule
\endfirsthead
\multicolumn{3}{@{}l}{\emph{(continued)}} \\
\toprule
Position & Chance of a loss & 95\% VaR \\
\midrule
\endhead
\midrule
\multicolumn{3}{r@{}}{\emph{(continued on the next page)}} \\
\endfoot
\bottomrule
\endlastfoot
Bond A alone & 4\% & 0, because 4\% is below 5\% \\*
Bond B alone & 4\% & 0, for the same reason \\
Sum of the two & & 0 \\*
Both bonds together & 7.84\% & 100, because 7.84\% is above 5\% \\
\end{longtable}
\addtocounter{table}{-1}% uncaptioned longtable: keep table numbers unchanged
\endgroup

Held separately, the bonds report no risk. Held together, they report a VaR of 100, so diversifying has \emph{increased} the measured risk.

The example points to a deeper problem. Suppose a trader sells insurance-like contracts that earn a small fee on most days but lose \$100 million on 1 day in 200, a probability of 0.5\%. VaR and ES judge this position very differently:

To see why, imagine the trader's results over 10,000 trading days. On 50 of those days (0.5\%), the trader loses \$100 million. On the other 9,950 days, the trader earns a small fee and loses almost nothing. Now rank all 10,000 days from worst to best. The worst 1\% are the 100 worst days: the 50 days with the \$100 million loss come first, followed by 50 ordinary days.


VaR reports almost no risk, while ES reveals it. This blind spot is one reason regulators replaced VaR with ES. It is the same trap as the put-selling trader of Chapter~\ref{ch:01}, who looks like a genius until the rare disaster arrives.

\begingroup
\small
\setlength{\LTpre}{6pt}\setlength{\LTpost}{6pt}
\begin{longtable}{@{}>{\raggedright\arraybackslash}p{0.14\linewidth}>{\raggedright\arraybackslash}p{0.22\linewidth}>{\raggedright\arraybackslash}p{0.14\linewidth}>{\raggedright\arraybackslash}p{\dimexpr 0.9\linewidth-0.14\linewidth-0.22\linewidth-0.14\linewidth-6\tabcolsep\relax}@{}}
\toprule
Measure & What it reports & Result & Why \\
\midrule
\endfirsthead
\multicolumn{4}{@{}l}{\emph{(continued)}} \\
\toprule
Measure & What it reports & Result & Why \\
\midrule
\endhead
\midrule
\multicolumn{4}{r@{}}{\emph{(continued on the next page)}} \\
\endfoot
\bottomrule
\endlastfoot
99\% VaR & The loss on the 100th worst day & About \$0 & The 100th worst day is an ordinary day \\*
99\% ES & The average loss on the 100 worst days & \$50 million & Half of the 100 days lose \$100 million \\*
\end{longtable}
\addtocounter{table}{-1}% uncaptioned longtable: keep table numbers unchanged
\endgroup


\section{Six ways to compute VaR}
\label{sec:11-how}

Every VaR model must answer two questions: how volatile will tomorrow be, and how fat is the tail of losses? The six models tested in this chapter combine different answers.

\begingroup
\small
\setlength{\LTpre}{6pt}\setlength{\LTpost}{6pt}
\begin{longtable}{@{}>{\raggedright\arraybackslash}p{0.28\linewidth}>{\raggedright\arraybackslash}p{\dimexpr(0.9\linewidth-0.28\linewidth-4\tabcolsep)/2\relax}>{\raggedright\arraybackslash}p{\dimexpr(0.9\linewidth-0.28\linewidth-4\tabcolsep)/2\relax}@{}}
\toprule
Model & Volatility & Tail shape \\
\midrule
\endfirsthead
\multicolumn{3}{@{}l}{\emph{(continued)}} \\
\toprule
Model & Volatility & Tail shape \\
\midrule
\endhead
\midrule
\multicolumn{3}{r@{}}{\emph{(continued on the next page)}} \\
\endfoot
\bottomrule
\endlastfoot
Normal, rolling window & Equal-weighted past window (here 250 days) & Normal \\*
Historical simulation (HS) & Implicit in the window & Empirical quantile of past returns \\
EWMA (RiskMetrics) & Exponentially weighted, $\lambda = 0.94$ & Normal \\
GARCH & Estimated, mean-reverting & Normal \\
GJR-GARCH with Student-$t$ shocks & Estimated, with leverage effect & Student-$t$, estimated \\*
Filtered historical simulation (FHS) & GJR-GARCH & Empirical quantile of past \emph{standardized} returns \\
\end{longtable}
\addtocounter{table}{-1}% uncaptioned longtable: keep table numbers unchanged
\endgroup

\term{Historical simulation}\index{Value at Risk (VaR)!historical simulation|textbf} ranks the returns of a past window from worst to best and reads off the VaR directly. With 500 days of history, for example, the 99\% VaR is the loss between the 5th and 6th worst days. It is the most popular method at banks (P\'erignon and Smith 2010), because it assumes no distribution and handles options by revaluing the portfolio under each past day's market moves. Its weakness is that it reacts slowly. A 250-day window gives a crash eleven months ago the same weight as yesterday's, so when volatility jumps, the VaR rises only as new large losses enter the window, and it drops abruptly when they leave.

\textbf{EWMA and GARCH}\index{EWMA|textbf} react to volatility within days, which is Chapter~\ref{ch:07}'s central finding. With normal shocks, however, they understate the tail. Even after volatility clustering is removed, daily returns have fatter tails than the normal distribution, so a normal 99\% VaR built on a good volatility forecast\index{volatility!forecasting} is still breached too often.

\term{Filtered historical simulation}\index{Value at Risk (VaR)!filtered historical simulation|textbf} (Barone-Adesi, Giannopoulos and Vosper 1999; Hull and White 1998) combines the strengths of both approaches. It first divides each past return by that day's volatility forecast, which puts calm and turbulent days on the same scale; these rescaled returns are called standardized returns. It then reads the VaR from the worst of these rescaled returns, as historical simulation does, which captures the fat tails\index{fat tails} without assuming a distribution. Finally, it multiplies by today's volatility forecast, so the VaR reacts quickly. The GJR\index{GARCH!GJR}-$t$ model follows the same idea but uses a Student-$t$ distribution, a bell-shaped curve with fatter tails than the normal, for the tail.


\section{The claim}
\label{sec:11-claim}

\begin{callout}
\textbf{Claim.} A well-built 99\% one-day VaR model is breached on about 1\% of trading days, and its breaches are independent over time. Breaches neither cluster in turbulent periods nor become more likely after a breach.
\end{callout}

This claim is implicit in every VaR number a bank reports and in the Basel rules that used VaR to set capital. Christoffersen (1998) called its two parts \term{unconditional coverage}\index{coverage test!unconditional|textbf} and \term{independence}. Together they form \term{conditional coverage}\index{coverage test!conditional|textbf}: given everything known the day before, the chance of a breach is exactly 1\%. Four points need to be settled.

\begin{itemize}
  \item \textbf{Which models.} The claim is that some models, built on the evidence of Chapter~\ref{ch:07}, achieve conditional coverage. We test six, from the simplest to filtered historical simulation.
  \item \textbf{Which portfolio.} The U.S. stock market, a single liquid position with daily data back to 1926. Bank portfolios are more complex, and their own backtests (Section~\ref{sec:11-evidence}) are the stronger test of practice.
  \item \textbf{Which confidence level.} The claim is stated at the Basel level of 99\%, but a model that fails at 99.9\% is of limited use for the tail events that matter.
  \item \textbf{What would count against it.} A breach rate significantly different from 1\%, or significantly clustered breaches. The claim weakens if every model fails in crises.
\end{itemize}

\section{The original evidence}
\label{sec:11-evidence}

\textbf{The Basel rules.}\index{Basel framework} From 1996, banks could set market-risk capital with their own models, at three times their 99\% ten-day VaR, with a higher multiplier if a 250-day backtest found too many breaches (the \textbf{traffic light}\index{Basel traffic light|textbf} of Section~\ref{sec:11-backtest}). A failed backtest costs a bank money.

\textbf{One year is too short.} Kupiec (1995) showed that a 99\% VaR should be breached only about 2.5 times in 250 days, so even a badly wrong model is hard to reject in one year.

\textbf{Banks' own models.} Berkowitz and O'Brien (2002) found that six large U.S. banks' VaRs were too cautious on average, but the breaches that did occur came together in the turmoil of autumn 1998. A simple GARCH model did about as well and reacted faster.

\textbf{Which models work best.} Comparisons of many models agree that a GARCH-type volatility forecast with a fat-tailed distribution performs best, and that plain historical simulation and normal models with slow-moving volatility perform worst (Kuester, Mittnik and Paolella 2006).

\section{Backtesting VaR and expected shortfall}
\label{sec:11-backtest}

A backtest compares each day's VaR forecast, made the day before, with that day's actual loss. A loss larger than the VaR is a \textbf{breach}\index{VaR exceedance|textbf}, or exception. Recording 1 for a breach and 0 otherwise turns a year of trading into a string of 0s and 1s, which every test in this section uses.

\textbf{Test 1: the right number of breaches (Kupiec 1995).}\index{coverage test!unconditional|textbf} A 99\% VaR should be breached about 2.5 times in 250 days. Kupiec's test asks whether the actual count is too far from this to be bad luck.

\begin{technote}
\textbf{Technical note: the Kupiec statistic.} With $n$ days, $x$ breaches and a promised breach rate $\alpha$,
\begin{equation}
  LR_{uc} = -2\ln\!\left[(1 - \alpha)^{n-x}\,\alpha^{x}\right] + 2\ln\!\left[\left(1 - \tfrac{x}{n}\right)^{n-x}\left(\tfrac{x}{n}\right)^{x}\right]
  \label{eq:11-kupiec}
\end{equation}
compares how likely the breaches are at the promised rate and at the observed rate. If the model is correct (the null hypothesis\index{hypothesis testing!null hypothesis}), $LR_{uc}$ follows a $\chi^2$ distribution with one degree of freedom, so a value above 3.84 rejects it at the 5\% level. The test rejects too many breaches, which understates risk, and too few, which wastes capital.
\end{technote}

\textbf{Test 2: breaches spread out (Christoffersen 1998).}\index{Christoffersen test} A model can have the right number of breaches all in one bad month. Christoffersen's test compares the chance of a breach after a normal day, $\pi_{01}$, with the chance after a breach, $\pi_{11}$; for a good 99\% model both are about 1\%. A $\pi_{11}$ of 5\% to 10\% means the model reacts too slowly to turbulence. This test is also $\chi^2(1)$, and adding it to Kupiec's gives the $\chi^2(2)$ \textbf{conditional-coverage test}\index{coverage test!conditional}.

\textbf{The Basel traffic light.}\index{Basel framework|textbf}\index{Basel traffic light} Regulators count breaches of a bank's 99\% one-day VaR over the last 250 days. More breaches mean more capital, as a multiple of VaR:

\begingroup
\small
\setlength{\LTpre}{6pt}\setlength{\LTpost}{6pt}
\begin{longtable}{@{}>{\raggedright\arraybackslash}p{0.12\linewidth}>{\raggedright\arraybackslash}p{0.16\linewidth}>{\raggedright\arraybackslash}p{0.22\linewidth}>{\raggedright\arraybackslash}p{\dimexpr 0.9\linewidth-0.12\linewidth-0.16\linewidth-0.22\linewidth-6\tabcolsep\relax}@{}}
\toprule
Zone & Breaches & Capital multiplier & Chance that a correct model has this many or fewer \\
\midrule
\endfirsthead
\multicolumn{4}{@{}l}{\emph{(continued)}} \\
\toprule
Zone & Breaches & Capital multiplier & Chance that a correct model has this many or fewer \\
\midrule
\endhead
\midrule
\multicolumn{4}{r@{}}{\emph{(continued on the next page)}} \\
\endfoot
\bottomrule
\endlastfoot
Green & 0--4 & 3.0 & 89.2\% \\*
Yellow & 5--9 & 3.40 to 3.85 & 99.99\% at 9 \\*
Red & 10 or more & 4.0 & n/a \\*
\end{longtable}
\addtocounter{table}{-1}% uncaptioned longtable: keep table numbers unchanged
\endgroup

% --- keep the traffic-light figure (fig:11-traffic) here, unchanged ---

Even a correct model lands in the yellow zone about one year in nine and almost never in the red. The zones balance two mistakes: punishing banks with good models and letting off banks with bad ones.

\textbf{Backtesting expected shortfall.}\index{expected shortfall (ES)|textbf} Since 2016, banks have set capital with 97.5\% ES. ES is an average size of loss, so no simple count can check it; statisticians say it is not \term{elicitable}\index{elicitability|textbf} (Gneiting 2011). The \term{exceedance-residual test} of McNeil and Frey (2000) checks whether losses on breach days match the ES forecast on average, and Part D of the lab applies it. Because ES is hard to test, Basel still backtests VaR at 97.5\% and 99\%.

\section{Lab~\thechapter: Replicating the Kupiec and Christoffersen backtests}
\label{sec:11-lab}

\textbf{Claim on trial\index{claims-on-trial method}:} losses exceed a 99\% VaR on about 1\% of trading days, in calm markets and in crises alike. The lab replicates the backtests of Kupiec (1995) and Christoffersen (1998) on daily U.S. market returns from the Kenneth French\index{Kenneth French data library} library, from 1930 to the present, about 24,000 trading days. It backtests the six models of Section~\ref{sec:11-how}, using only data available the day before and re-estimating the GARCH models yearly on the previous ten years. A single liquid stock-market position is the simplest setting for VaR, so if the models struggle here, they will not do better with a bank's options, credit and hard-to-value positions.

\textbf{Part A, the number of breaches.} Breaches of the 99\% VaR against the roughly 240 expected, with Kupiec's test, the worst loss as a multiple of VaR, and the share of years in the yellow or red zone (Lab Figures~\ref{ch:08}.1 and~\ref{ch:08}.2).

\textbf{Part B, the timing of breaches.} The chance of a breach after a breach, and Christoffersen's test (Lab Figure~\ref{ch:08}.3).

\textbf{Part C, robustness checks.} The backtests under three harder conditions (Section~\ref{sec:11-stress}).

\textbf{Part D, expected shortfall.} At 97.5\%, the average loss on breach days against the average ES forecast, and the McNeil--Frey test (Lab Figure~\ref{ch:08}.5).

\label{sec:11-stress}
\textbf{Robustness checks.} On October 19, 1987, the Dow fell 22.6\%, a loss no model built on the calm months before could have put inside a 99\% VaR, so Part C counts breaches in 1930--32, October 1987, autumn 2008 and spring 2020 and measures how quickly each model catches up. Deeper in the tail, a normal model that is slightly off at 95\% can be several times off at 99.9\% (Lab Figure~\ref{ch:08}.4), and with only about 24 breaches expected in 95 years the test is also weak there. Finally, no equal-weighted\index{index weighting!equal-weighted} historical-simulation window stops breaches from bunching, because none reacts to volatility as GARCH does.

\section{Where the evidence stands}
\label{sec:11-verdict}

The studies and records cited in this chapter point to a clear pattern: \textbf{VaR can measure risk reliably, but only with the right model, and not during a crisis.} The table summarizes the evidence as of November 2026.

\begingroup
\small
\setlength{\LTpre}{6pt}\setlength{\LTpost}{6pt}
\begin{longtable}{@{}>{\raggedright\arraybackslash}p{0.17\linewidth}>{\raggedright\arraybackslash}p{0.13\linewidth}>{\raggedright\arraybackslash}p{\dimexpr \linewidth-0.17\linewidth-0.13\linewidth-0.22\linewidth-6\tabcolsep\relax}>{\raggedright\arraybackslash}p{0.22\linewidth}@{}}
\toprule
Situation & Does the claim hold? & Why & Evidence \\
\midrule
\endfirsthead
\multicolumn{4}{@{}l}{\emph{(continued)}} \\
\toprule
Situation & Does the claim hold? & Why & Evidence \\
\midrule
\endhead
\midrule
\multicolumn{4}{r@{}}{\emph{(continued on the next page)}} \\
\endfoot
\bottomrule
\endlastfoot
The best models in normal times & \textbf{Yes} & GARCH-type models with fat tails are breached on close to 1\% of days, with little bunching. & Kuester, Mittnik and Paolella (2006); Berkowitz and O'Brien (2002) \\*
Spotting bad models & \textbf{Yes} & Backtests catch models with too many or bunched breaches. & Kupiec (1995); Christoffersen (1998) \\
The models most banks use & \textbf{No} & Historical simulation reacts slowly, so its breaches bunch when losses matter most. & P\'erignon and Smith (2010); Berkowitz and O'Brien (2002) \\
The first days of a crisis & \textbf{No} & No model built on past data foresees a crash. & Berkowitz and O'Brien (2002) on 1998; the UBS investigation on 2007--08; Federal Reserve (2021) on 2020 \\
Deep in the tail (99.9\%) & \textbf{No} & Fat tails make normal models fail more the further out they go, and too few breaches are expected to test well. & Fat tails (Chapter~\ref{ch:07}); Section~\ref{sec:11-stress} \\*
The size of the worst losses & \textbf{Not measured} & VaR\index{Value at Risk (VaR)!backtesting} shows where the bad days begin, not how bad they are. & Artzner et al.\ (1999); Basel's move to ES (2016) \\
\end{longtable}
\addtocounter{table}{-1}% uncaptioned longtable: keep table numbers unchanged
\endgroup

\section{What practitioners and investors can use}
\label{sec:11-investors}

Anyone measuring risk should use a volatility forecast that reacts quickly, such as EWMA or a GARCH-type model, allow for fat tails, report ES alongside VaR, backtest\index{backtest|textbf} all the time, and run scenario tests such as a repeat of October 1987, 2008 or 2022. A fast-reacting VaR has one side effect: it raises capital requirements just as markets fall, so banks sell together and can deepen the crisis, a problem called \term{procyclicality}. Basel 2.5 (2009) therefore added a \term{stressed VaR}\index{Value at Risk (VaR)!stressed|textbf}, measured over a year of severe stress. Investors reading a bank's or fund's risk reports should ask how many breaches there were last year and when, and how large the stressed VaR is compared with the reported VaR.

\textbf{What Failed: bank VaR in 2007--08.} By 2007, large banks measured trading risk mainly with VaR. UBS suffered the largest subprime losses of any European bank. The Swiss Federal Banking Commission found that UBS had judged its super-senior CDO tranches\index{tranches}, the supposedly safest slices of securities built from mortgage-related bonds, to be ``VaR neutral,'' so they ``were not reflected in key internal risk reports.'' By UBS's own account, these positions caused about half of its losses by the end of 2007. VaR was not calculated wrongly. It measured positions that looked safe in recent data and could not see the risk of a market that had never been stressed.

% Key terms of the chapter, in alphabetical order, for readers new to finance.
% A term defined in an earlier chapter points to the chapter, section and page.
\providecolor{FEMaroon}{RGB}{128,30,45}
\providecolor{FECream}{RGB}{253,248,238}
\providecolor{FEGold}{RGB}{196,150,60}
\makeatletter
\@ifundefined{keyterms}{%
  \newtcolorbox{keyterms}[1][Key terms in this chapter]{breakable, enhanced,
    colback=FECream, colframe=FEGold, boxrule=0.6pt, arc=2pt, leftrule=3pt,
    left=6pt, right=6pt, top=3pt, bottom=3pt, before skip=8pt, after skip=8pt,
    fonttitle=\bfseries, coltitle=white, colbacktitle=FEMaroon, title={#1}}%
}{}
\makeatother

\begin{keyterms}
\footnotesize
\begin{description}[leftmargin=0pt, itemsep=0pt, parsep=0pt, topsep=0pt, font=\normalfont\bfseries\color{FEMaroon}]
  \item[Basel Committee.] The Basel Committee on Banking Supervision is the international group of bank regulators that sets the rules for banks' capital.
  \item[Breach.] A breach is a day on which the loss is larger than the VaR forecast.
  \item[Christoffersen test.] The Christoffersen (1998) test checks whether VaR breaches arrive independently or in clusters.
  \item[Conditional coverage.] Conditional coverage combines unconditional coverage and independence, so that the chance of a breach is exactly as promised given everything known the day before.
  \item[Elicitable.] A risk measure is elicitable when some scoring rule can rank forecasts of it against outcomes, which VaR is and expected shortfall alone is not.
  \item[Exceedance-residual test.] The exceedance-residual test checks whether the losses on breach days match the expected-shortfall forecast on average.
  \item[Expected shortfall (ES).] Expected shortfall is the average loss on the worst days, those beyond the VaR.
  \item[Filtered historical simulation.] Filtered historical simulation rescales past returns by a volatility model before ranking them, combining the strengths of historical simulation and GARCH.
  \item[Historical simulation.] Historical simulation ranks the returns of a past window from worst to best and reads the VaR directly from them.
  \item[Independence of breaches.] Independence means that a VaR breach today does not make a breach tomorrow more likely, so breaches do not come in clusters.
  \item[Kupiec test.] The Kupiec (1995) test checks whether VaR breaches occur as often as the model promises.
  \item[Procyclicality.] Procyclicality is the way risk models that react to recent volatility force many banks to cut risk at the same time in a crisis, deepening it.
  \item[Quantile.] A quantile, or percentile, is the value below which a given share of outcomes falls, so the 1st percentile is exceeded on the downside only 1\% of the time.
  \item[Stressed VaR.] Stressed VaR is a VaR measured over a year of severe market stress, which regulators added to the capital rules in 2009.
  \item[Subadditive.] A risk measure is subadditive when the risk of a combined portfolio is never more than the sum of the risks of its parts, so it always rewards diversification.
  \item[Unconditional coverage.] Unconditional coverage means that a VaR model is breached as often as promised on average, for example on 1\% of days.
  \item[Value at Risk (VaR).] See the introduction to Chapter~\ref{ch:07}, page~\pageref{kt:value-at-risk}.
\end{description}
\end{keyterms}
